\documentclass[conference]{IEEEtran}
\usepackage{cite}
\usepackage{amsmath,amssymb,amsfonts}
\usepackage{booktabs}
\usepackage{url}
\usepackage{hyperref}

\begin{document}

% ─── TITLE ───────────────────────────────────────────────────────────────────
\title{Predicting Employee Turnover: A Statistical Analysis\\of Workforce Attrition Factors}

\author{
  \IEEEauthorblockN{Brenda Murillo}
  \IEEEauthorblockA{
    \textit{Dept. of Computer and Electrical Engineering and Computer Science}\\
    \textit{California State University, Bakersfield}\\
    Bakersfield, CA, USA\\
    bmurillo3@csub.edu
  }
}

\maketitle

% ─── ABSTRACT ────────────────────────────────────────────────────────────────
\begin{abstract}
Employee turnover represents one of the most persistent and costly challenges
facing organizations today, yet most retention strategies rely on intuition
rather than data. This paper presents a statistical analysis of workforce
attrition using an HR dataset of 14,999 employee records drawn from Kaggle.
We apply descriptive statistics, relational database querying (SQLite),
correlation analysis, and visual analysis to identify the factors most
strongly associated with employee departure. Our findings reveal an overall
turnover rate of 23.81\%, with low-salary employees leaving at nearly
30\%---more than four times the rate of high-salary employees. Job
satisfaction emerges as the single strongest predictor, with a Pearson
correlation of $r = -0.388$ with turnover; employees who left averaged a
satisfaction score of 0.44 compared to 0.67 for those who stayed. We further
identify 615 current employees who match a high-risk churn profile (low
satisfaction, low salary, no promotion in five years). These findings provide
HR departments with actionable, data-driven criteria for early intervention
and targeted retention strategies.
\end{abstract}

\begin{IEEEkeywords}
employee turnover, workforce attrition, HR analytics, churn prediction,
descriptive statistics, SQLite, data analysis
\end{IEEEkeywords}

% ─── I. INTRODUCTION ─────────────────────────────────────────────────────────
\section{Introduction}

The cost of replacing an employee is widely estimated at between 50\% and
200\% of that employee's annual salary, accounting for recruiting, onboarding,
lost productivity, and institutional knowledge transfer~\cite{griffeth2000}.
Yet despite this substantial burden, many organizations respond to turnover
reactively---conducting exit interviews only after employees have already
resigned, and implementing changes too late to retain the individuals who
prompted them.

The central challenge is one of prediction: given observable workplace
factors, can an organization identify which employees are most likely to leave
before they submit their resignation? This problem is well-suited to
data-driven analysis. Modern HR systems routinely capture variables such as
satisfaction surveys, performance evaluations, compensation tiers, workload
metrics, and promotion histories---precisely the kinds of signals that may
indicate attrition risk.

This paper addresses that challenge by applying statistical analysis to an HR
dataset containing records for 14,999 employees. We pursue four core research
questions:
\begin{enumerate}
  \item How strongly does salary level predict employee turnover?
  \item What is the relationship between job satisfaction and the likelihood
        of departure?
  \item Does workload (measured in monthly hours and project count) influence
        turnover?
  \item Can we identify current employees at high churn risk using observable
        characteristics?
\end{enumerate}

Our approach is deliberately grounded in accessible statistical
methods---descriptive statistics, group comparisons, correlation analysis, and
SQL querying against a normalized relational database---rather than black-box
machine learning models. This makes our findings interpretable and actionable
for HR practitioners who may not have data science expertise.

The remainder of this paper is organized as follows. Section~\ref{sec:related}
reviews related work. Section~\ref{sec:dataset} describes the dataset.
Section~\ref{sec:method} details our methodology. Section~\ref{sec:results}
presents results. Section~\ref{sec:discussion} discusses implications.
Section~\ref{sec:conclusion} concludes with limitations and future directions.

% ─── II. RELATED WORK ────────────────────────────────────────────────────────
\section{Related Work}
\label{sec:related}

Employee attrition prediction has been an active area of research in both
human resources management and data science. Early work by Mobley~\cite{mobley1977}
established a foundational behavioral model of turnover, identifying job
dissatisfaction as the primary precursor to resignation---a finding
consistently replicated in subsequent literature. Griffeth, Hom, and
Gaertner~\cite{griffeth2000} later confirmed through meta-analysis that
satisfaction, pay, and workload are the most reliable predictors of departure
across industries and time periods.

More recent computational approaches have applied supervised machine learning
to HR datasets. Punnoose and Ajit~\cite{punnoose2016} applied random forests
and logistic regression to the same Kaggle HR dataset used in this study,
achieving classification accuracies above 97\%. Similarly, Alduayj and
Rajpoot~\cite{alduayj2018} compared decision trees, support vector machines,
and neural networks for attrition prediction, finding that tree-based models
consistently outperformed others on HR tabular data. Zhao et al.~\cite{zhao2018}
demonstrated that ensemble methods further improve predictive performance but
at significant cost to interpretability.

While the predictive accuracy of machine learning models is high, a
recognized limitation is interpretability~\cite{doshi2017}. HR managers need
to understand \textit{why} a model flags an employee as high-risk, not just
that it does. Our work is positioned in the tradition of transparent
statistical approaches that make decision criteria visible and reproducible
by any analyst.

% ─── III. DATASET ────────────────────────────────────────────────────────────
\section{Dataset}
\label{sec:dataset}

\subsection{Source and Scope}

The dataset is the HR Employee Retention dataset publicly available on
Kaggle~\cite{kaggle}, compiled to simulate realistic human resources records.
It contains \textbf{14,999 employee records}, each representing a unique
individual. The outcome variable (\texttt{left}) indicates whether the
employee departed the organization (1) or remained (0).

\subsection{Variables}

The dataset includes ten variables spanning four conceptual domains, shown
in Table~\ref{tab:variables}.

\begin{table}[h]
\caption{Dataset Variables}
\label{tab:variables}
\centering
\begin{tabular}{llll}
\toprule
\textbf{Domain} & \textbf{Variable} & \textbf{Type} & \textbf{Description} \\
\midrule
Satisfaction  & \texttt{satisfaction\_level}   & Numeric [0,1] & Self-reported satisfaction \\
Performance   & \texttt{last\_evaluation}       & Numeric [0,1] & Last performance review \\
Workload      & \texttt{number\_project}        & Integer       & Projects assigned \\
Workload      & \texttt{average\_monthly\_hours}& Integer       & Avg.\ hours/month \\
Tenure        & \texttt{time\_spend\_company}   & Integer       & Years at company \\
Safety        & \texttt{work\_accident}         & Boolean       & Had a work accident \\
Career        & \texttt{promotion\_last\_5years}& Boolean       & Promoted in 5 yrs \\
Compensation  & \texttt{salary}                 & Categorical   & Low / medium / high \\
Outcome       & \texttt{left}                   & Boolean       & Left the company \\
\bottomrule
\end{tabular}
\end{table}

\subsection{Data Quality}

A missing value audit revealed no null entries in any column. Outlier
detection using the interquartile range (IQR) method identified a small number
of anomalous values in \texttt{average\_monthly\_hours} (values above 310
hours/month), which were retained as plausible extreme-workload cases. Schema
constraints enforced data integrity: \texttt{CHECK} clauses validated that
satisfaction and evaluation scores fall within [0,1], that hours and tenure
are non-negative, and that foreign key relationships across all four tables
are consistent. A row-count reconciliation confirmed all 14,999 records are
present and accounted for after normalization.

\subsection{Database Schema}

To support efficient querying, the flat CSV was normalized into a four-table
relational schema in SQLite~\cite{codd1970}:
\begin{itemize}
  \item \texttt{salary\_levels} --- lookup table mapping salary tier labels to IDs
  \item \texttt{employees} --- core identity table (emp\_id, salary\_id, tenure)
  \item \texttt{performance\_records} --- satisfaction, evaluation, project count,
        monthly hours
  \item \texttt{employment\_events} --- work accidents, promotion history, and
        turnover outcome
\end{itemize}
This normalized design eliminates redundancy, enforces referential integrity
via foreign keys, and allows the turnover outcome to be joined against any
combination of employee attributes in a single SQL query.

% ─── IV. METHODOLOGY ─────────────────────────────────────────────────────────
\section{Methodology}
\label{sec:method}

\subsection{Descriptive Statistics}

For all numeric variables, we computed mean, median, mode, standard deviation,
variance, range, quartiles (Q1, Q3), and the interquartile range. These
statistics were calculated separately for employees who left and those who
stayed, enabling direct comparison of distributional differences across
turnover groups.

\subsection{Group Comparison}

Turnover rates were computed for each level of the categorical variables
(salary tier, promotion status, work accident occurrence). For numeric
variables, we segmented employees into bins and computed turnover rates within
each bin to identify non-linear relationships.

\subsection{Correlation Analysis}

Pearson correlation coefficients were computed between each numeric feature
and the binary \texttt{left} outcome, producing a ranked list of predictors by
association strength.

\subsection{SQL-Based Analysis}

All group-level queries were executed against the normalized SQLite database.
Key queries included: overall turnover rate (\texttt{COUNT} with
\texttt{GROUP BY left\_company}); turnover rate by salary tier (three-table
join across \texttt{employees}, \texttt{salary\_levels},
\texttt{employment\_events}); average satisfaction and monthly hours by
turnover status; and high-risk employee identification using compound
\texttt{WHERE} filters.

\subsection{Risk Profiling}

A high-risk profile was defined using three simultaneous criteria:
satisfaction score below 0.4, salary tier of ``low'', and no promotion in
the last five years, constrained to employees still at the company. This
threshold was set based on the distributional boundary at which turnover
rates were observed to increase sharply.

% ─── V. RESULTS ──────────────────────────────────────────────────────────────
\section{Results}
\label{sec:results}

\subsection{Overall Turnover Rate}

Of 14,999 employees, 3,571 left the organization, yielding an \textbf{overall
turnover rate of 23.81\%}. This substantially exceeds the commonly cited U.S.
industry average of approximately 15\%~\cite{bls2023}.

\subsection{Turnover by Salary Tier}

Salary tier is one of the strongest single-variable predictors of turnover.
Table~\ref{tab:salary} shows turnover rates disaggregated by compensation
group. Low-salary employees leave at more than \textbf{four times} the rate
of high-salary employees.

\begin{table}[h]
\caption{Turnover Rate by Salary Tier}
\label{tab:salary}
\centering
\begin{tabular}{lrrr}
\toprule
\textbf{Salary Tier} & \textbf{Total} & \textbf{Left} & \textbf{Rate} \\
\midrule
Low    & 7,316 & 2,172 & \textbf{29.69\%} \\
Medium & 6,446 & 1,317 & 20.43\% \\
High   & 1,237 &    82 & \textbf{6.63\%} \\
\bottomrule
\end{tabular}
\end{table}

\subsection{Job Satisfaction and Turnover}

Job satisfaction shows the most pronounced difference between retained and
departed employees across all variables (Table~\ref{tab:avgs}). It also
produces the strongest Pearson correlation with turnover outcome
($r = -0.388$).

\begin{table}[h]
\caption{Average Feature Values by Turnover Status}
\label{tab:avgs}
\centering
\begin{tabular}{lcc}
\toprule
\textbf{Feature} & \textbf{Stayed} & \textbf{Left} \\
\midrule
Satisfaction level      & \textbf{0.667} & \textbf{0.440} \\
Last evaluation score   & 0.715 & 0.718 \\
Number of projects      & 3.79  & 3.86  \\
Avg.\ monthly hours     & 199.1 & 207.4 \\
Years at company        & 3.38  & 3.88  \\
\bottomrule
\end{tabular}
\end{table}

Satisfaction differs by 0.227 points between groups---a 34\% relative gap.
By contrast, evaluation scores are nearly identical (0.715 vs.\ 0.718),
confirming that turnover is a satisfaction problem, not a performance problem.

\subsection{Feature Correlation Analysis}

Table~\ref{tab:corr} presents the Pearson correlation of each feature with
the \texttt{left} outcome, ranked by absolute magnitude.

\begin{table}[h]
\caption{Feature Correlation with Turnover ($\mathtt{left} = 1$)}
\label{tab:corr}
\centering
\begin{tabular}{lc}
\toprule
\textbf{Feature} & \textbf{Pearson $r$} \\
\midrule
\texttt{satisfaction\_level}    & $-0.388$ \\
\texttt{work\_accident}         & $-0.155$ \\
\texttt{time\_spend\_company}   & $+0.145$ \\
\texttt{average\_monthly\_hours}& $+0.071$ \\
\texttt{promotion\_last\_5years}& $-0.062$ \\
\texttt{number\_project}        & $+0.024$ \\
\texttt{last\_evaluation}       & $+0.007$ \\
\bottomrule
\end{tabular}
\end{table}

Satisfaction is the strongest single predictor by a wide margin.
\texttt{last\_evaluation} shows virtually no correlation ($r = +0.007$),
confirming that performance evaluation does not distinguish leavers from
stayers.

\subsection{Workload and Turnover}

Employees who left worked an average of \textbf{207.4 hours per month}
compared to 199.1 for those who stayed---a difference of 8.3 hours, or
roughly two extra workdays per month. The correlation of
\texttt{average\_monthly\_hours} with turnover ($r = +0.071$) suggests
overwork is a contributing factor, though it acts as a compounding stressor
alongside low satisfaction and low pay rather than an independent driver.

\subsection{High-Risk Employee Profiling}

Applying the compound risk profile, the analysis identified \textbf{615
current employees} who match all three criteria. Table~\ref{tab:risk} shows
the ten highest-risk individuals by lowest satisfaction score.

\begin{table}[h]
\caption{Sample High-Risk Employees (Top 10 by Lowest Satisfaction)}
\label{tab:risk}
\centering
\begin{tabular}{cccr}
\toprule
\textbf{Emp.\ ID} & \textbf{Satisfaction} & \textbf{Hrs/Mo} & \textbf{Tenure} \\
\midrule
2087  & 0.12 & 244 & 5 \\
3129  & 0.12 & 257 & 6 \\
4026  & 0.12 & 241 & 2 \\
4684  & 0.12 & 276 & 4 \\
5336  & 0.12 & 166 & 3 \\
7772  & 0.12 & 161 & 4 \\
8745  & 0.12 & 287 & 4 \\
9283  & 0.12 & 200 & 3 \\
11069 & 0.12 & 191 & 5 \\
13280 & 0.12 & 191 & 5 \\
\bottomrule
\end{tabular}
\end{table}

% ─── VI. DISCUSSION ──────────────────────────────────────────────────────────
\section{Discussion}
\label{sec:discussion}

\subsection{Compensation is the Dominant Structural Factor}

The four-to-one ratio in turnover rates between low and high salary tiers is
the single most actionable finding in this study. It suggests that
compensation structure creates a systemic attrition risk that cannot be
resolved through engagement programs or managerial coaching alone.
Organizations with large proportions of low-salary employees should treat
compensation review as a primary retention strategy~\cite{griffeth2000}.

\subsection{Satisfaction Predicts Turnover Independently of Performance}

The near-identical evaluation scores between leavers and stayers (0.718
vs.\ 0.715), combined with the near-zero correlation of \texttt{last\_evaluation}
with turnover ($r = +0.007$), challenges the common assumption that turnover
is driven by underperformers. The data shows the opposite: departed employees
were equally well-evaluated as retained ones. This replicates Mobley's
foundational finding~\cite{mobley1977} and underscores that retention efforts
should be framed as a satisfaction and engagement problem, not a performance
management problem.

\subsection{Workload Acts as a Compounding Risk Factor}

The highest-risk employees in the profiling analysis combine low satisfaction
with above-average hours, suggesting that overwork amplifies the effect of low
satisfaction rather than driving attrition independently. HR interventions
should target the intersection of these factors, not either factor in
isolation~\cite{hom2017}.

\subsection{Early Intervention Opportunity}

The identification of 615 high-risk current employees is the most practically
valuable output of this analysis. Unlike post-departure analyses, this
profiling approach is actionable before attrition occurs. HR departments can
operationalize it as a recurring SQL query to generate a prioritized list for
manager outreach, salary review, or promotion consideration.

\subsection{Limitations}

This study has several limitations. First, the dataset is simulated~\cite{kaggle}
and does not represent a specific organization, which limits the
generalizability of specific threshold values. Second, our analysis is
descriptive rather than causal: we identify associations, not causal effects.
Third, the dataset lacks department or role information, which likely moderates
many of the relationships observed. Fourth, class imbalance (approximately
24\% leavers) should be addressed before extending this framework to predictive
modeling.

% ─── VII. CONCLUSION ─────────────────────────────────────────────────────────
\section{Conclusion}
\label{sec:conclusion}

This paper analyzed an HR dataset of 14,999 employees to identify the factors
most strongly associated with workforce attrition. Salary tier and job
satisfaction are the two dominant predictors of turnover---with satisfaction
producing the strongest correlation ($r = -0.388$)---while performance
evaluation is essentially uncorrelated with departure ($r = +0.007$).
Critically, departed employees were equally well-evaluated as retained
employees, confirming that turnover is a satisfaction problem, not a
performance problem.

By normalizing a flat HR CSV into a constraint-validated SQLite schema and
operationalizing findings as reusable SQL queries, we demonstrate that
organizations can generate actionable, prioritized intervention lists from
existing HR data without requiring machine learning infrastructure. The 615
high-risk employees identified here represent an immediately addressable
retention opportunity.

Future work should apply this framework to real organizational data with
department-level granularity, incorporate time-series analysis to model
attrition risk over an employee's tenure, and explore causal inference
methods (e.g., propensity score matching) to move beyond association toward
actionable causal claims about the effect of salary increases or promotion on
retention probability.

% ─── ACKNOWLEDGMENT ──────────────────────────────────────────────────────────
\section*{Acknowledgment}

The author thanks the faculty of the Department of Computer and Electrical
Engineering and Computer Science at California State University, Bakersfield
for guidance on this project.

% ─── REFERENCES ──────────────────────────────────────────────────────────────
\begin{thebibliography}{10}

\bibitem{mobley1977}
W.~H. Mobley, ``Intermediate linkages in the relationship between job
satisfaction and employee turnover,'' \textit{J. Appl. Psychol.}, vol.~62,
no.~2, pp.~237--240, 1977.

\bibitem{griffeth2000}
R.~W. Griffeth, P.~W. Hom, and S.~Gaertner, ``A meta-analysis of antecedents
and correlates of employee turnover: Update, moderator tests, and research
implications for the next millennium,'' \textit{J. Manag.}, vol.~26, no.~3,
pp.~463--488, 2000.

\bibitem{hom2017}
P.~W. Hom, T.~W. Lee, J.~D. Shaw, and J.~P. Hausknecht, ``One hundred years
of employee turnover theory and research,'' \textit{J. Appl. Psychol.},
vol.~102, no.~3, pp.~530--545, 2017.

\bibitem{bls2023}
Bureau of Labor Statistics, ``Job Openings and Labor Turnover Summary,''
U.S. Department of Labor, 2023. [Online]. Available:
\url{https://www.bls.gov/news.release/jolts.nr0.htm}

\bibitem{punnoose2016}
R.~Punnoose and P.~Ajit, ``Prediction of employee turnover in organizations
using machine learning algorithms,'' \textit{Int. J. Adv. Res. Artif.
Intell.}, vol.~5, no.~9, pp.~22--26, 2016.

\bibitem{zhao2018}
Y.~Zhao, M.~K. Hryniewicki, F.~Cheng, B.~Fu, and X.~Zhu, ``Employee turnover
prediction with machine learning: A reliable approach,'' in \textit{Proc.
Intelligent Systems Conf. (IntelliSys)}, 2018, pp.~737--758.

\bibitem{alduayj2018}
S.~S. Alduayj and K.~Rajpoot, ``Predicting employee attrition using machine
learning,'' in \textit{Proc. Int. Conf. Innovations in Information Technology
(IIT)}, 2018, pp.~93--98.

\bibitem{doshi2017}
F.~Doshi-Velez and B.~Kim, ``Towards a rigorous science of interpretable
machine learning,'' \textit{arXiv preprint arXiv:1702.08608}, 2017.

\bibitem{kaggle}
HR Analytics / Employee Retention Dataset, Kaggle. [Online]. Available:
\url{https://www.kaggle.com/}

\bibitem{codd1970}
E.~F. Codd, ``A relational model of data for large shared data banks,''
\textit{Commun. ACM}, vol.~13, no.~6, pp.~377--387, 1970.

\end{thebibliography}

\end{document}
